% Part: turing-machines
% Chapter: undecidability
% Section: representing-tms

\documentclass[../../../include/open-logic-section]{subfiles}

\begin{document}

% SEGMENT OLP-0270-S01

\olfileid{tur}{und}{rep}
\olsection{டியூரிங் பொறிகளைக் குறிப்பிடுதல்}

\begin{explain}
டியூரிங் பொறிகளையும் அவற்றின் நடத்தையையும் முதல்தர தருக்கத்தின்
!!a{sentence} ஆல் குறிப்பிடுவதற்கு, பொருத்தமான ஒரு மொழியை வரையறுக்க
வேண்டும். அம்மொழி இரு பகுதிகளைக் கொண்டது: பொறியின் அமைவுநிலைகளை
விவரிப்பதற்கான !!{predicate}s; செயலாக்கப் படிகளுக்கும்
(``கணங்களுக்கும்'') நாடாவின் இருப்பிடங்களுக்கும் எண்களிடுவதற்கான
கோவைகள்.

இரண்டு வகையான !!{predicate}s ஐ அறிமுகப்படுத்துகிறோம்; இரண்டும்
2-இடம் கொண்டவை: ஒவ்வொரு நிலை~$q$ க்கும் !!a{predicate}~$\Obj Q_q$,
ஒவ்வொரு நாடாக் குறி~$\sigma$ க்கும் !!a{predicate}~$\Obj S_\sigma$.
முந்தையவை~$M$ இன் நிலையையும் அதன் நாடாத் தலையின் இருப்பிடத்தையும்
விவரிக்க அனுமதிக்கின்றன; பிந்தையவை நாடாவின் உள்ளடக்கத்தை விவரிக்க
அனுமதிக்கின்றன.

நாடாத் தலையின் இருப்பிடங்களையும் செயலாக்கப்பட்ட படிகளின் எண்ணிக்கையையும்
வெளிப்படுத்த, எண்களை வெளிப்படுத்தும் ஒரு வழி தேவை. இதற்கு
!!a{constant}~$\Obj 0$ ஐயும், அடுத்தெண் சார்பான $\prime$ என்ற $1$-இடச்
சார்பையும் பயன்படுத்துகிறோம். மரபுப்படி, அது தனது மாறிக்குப்
\emph{பின்னால்} எழுதப்படுகிறது (அடைப்புக்குறிகளை நாம் விட்டுவிடுகிறோம்).

ஒவ்வோர் எண் $n$ க்கும்,~$\Lang L_M$ இல் அதைக் குறிப்பிடும்
~$n$ க்கான \emph{எண்ணுருவான} ஒரு நியமச் சொல்~$\num{n}$ உள்ளது.
$\num{0}$ என்பது $\Obj 0$, $\num{1}$ என்பது $\Obj 0'$, $\num{2}$
என்பது $\Obj 0''$, இவ்வாறாகத் தொடரும். மேலும் முறைசார்ந்து:
\begin{align*}
\num{0} & = \Obj 0 \\
\num{n+1} &= \num{n}'
\end{align*}
சொல் $\num{0}$, அதாவது $\Obj 0$, நாடாவின் இடக்கோடி இருப்பிடத்துக்கும்
முதல் செயலாக்கப் படிக்கு முந்தைய நேரத்துக்கும் (தொடக்க அமைவுநிலைக்கும்)
பெயரிடுகிறது. சொல் $\num{1}$, அதாவது $\Obj 0'$, இடக்கோடிக் கட்டத்துக்கு
வலப்புறமுள்ள கட்டத்துக்கும் முதல் செயலாக்கப் படிக்குப் பின்னான
நேரத்துக்கும் பெயரிடுகிறது; இவ்வாறாகத் தொடரும்.

நாடாவின் இருப்பிடங்களின் வரிசையையும் (அப்போது அதற்கு ``இதற்கு
இடப்புறத்தில்'' என்று பொருள்) செயலாக்கப் படிகளின் வரிசையையும்
(அப்போது அதற்கு ``இதற்கு முன்'' என்று பொருள்) வெளிப்படுத்த
!!a{predicate}~$<$ ஐயும் அறிமுகப்படுத்துகிறோம்.

மொழி அமைந்ததும், $!T(M, w)$ இன் ``அடிகோள்களை'', அதாவது
ஒன்றாக எடுத்துக்கொள்ளும்போது உள்ளீடு~$w$ இல் இயக்கப்படும்~$M$ இன்
நடத்தையை விவரிக்கும் !!{sentence}s ஐப் பட்டியலிடுகிறோம்.
$\Obj 0$, $\prime$, $<$ ஆகியவற்றின்மீது நிபந்தனைகளை விதிக்கும்
!!{sentence}s, உள்ளீட்டு அமைவுநிலையை விவரிக்கும் !!{sentence}s,
ஒரு குறிப்பிட்ட கட்டளையைச் செயலாக்கிய பிறகு~$M$ இன் அமைவுநிலை
என்னவென்று விவரிக்கும் !!{sentence}s ஆகியவை இருக்கும்.
\end{explain}

\begin{defn}
  \ollabel{defn:tm-descr}
ஒரு டியூரிங் பொறி $M = \tuple{Q, \Sigma, q_0, \delta}$
கொடுக்கப்பட்டால், மொழி~$\Lang L_M$ பின்வருவனவற்றைக் கொண்டது:
\begin{enumerate}
\item ஒவ்வொரு நிலை~$q \in Q$ க்கும் இரண்டு இடம் கொண்ட
!!{predicate} $\Obj Q_q(x, y)$. உள்ளுணர்வுப்படி,
$\Obj Q_q(\num{m}, \num{n})$ என்பது ``$n$ படிகளுக்குப் பிறகு,
~$M$ நிலை~$q$ இல் இருந்து $m$-ஆம் கட்டத்தை வருடுகிறது'' என்பதை
வெளிப்படுத்துகிறது.
\item ஒவ்வொரு குறி~$\sigma\in \Sigma$ க்கும் இரண்டு இடம் கொண்ட
!!{predicate} $\Obj S_\sigma(x, y)$. உள்ளுணர்வுப்படி,
$\Obj S_\sigma(\num{m}, \num{n})$ என்பது ``$n$ படிகளுக்குப் பிறகு,
$m$-ஆம் கட்டம் குறி~$\sigma$ ஐக் கொண்டுள்ளது'' என்பதை வெளிப்படுத்துகிறது.
\item !!a{constant} $\Obj 0$
\item ஓரிடம் கொண்ட !!a{function} $\prime$
\item இரு இடம் கொண்ட !!a{predicate} $<$
\end{enumerate}
\end{defn}

உள்ளீடு $w = \sigma_{i_1}\dots\sigma_{i_k}$ இல் டியூரிங் பொறி~$M$
செயல்படுவதை விவரிக்கும் !!{sentence}s பின்வருவன:
\begin{enumerate}
\item எண்களையும்~$<$ ஐயும் விவரிக்கும் அடிகோள்கள்:
\begin{enumerate}
\item ஒவ்வோர் எண்ணும் அதன் அடுத்தெண்ணைவிடச் சிறியது என்று கூறும்
!!^a{sentence}:
\[
\lforall[x][x < x']
\]
\item $<$ கடப்புநிலையுடையது என்பதை உறுதிசெய்யும் !!^a{sentence}:
\[
\lforall[x][\lforall[y][\lforall[z][
      ((x < y \land y < z) \lif x < z)]]]
\]
\end{enumerate}
\item உள்ளீட்டு அமைவுநிலையை விவரிக்கும் அடிகோள்கள்:
\begin{enumerate}
\item $0$ படிகளுக்குப் பிறகு---பொறி தொடங்குவதற்கு முன்---~$M$
  தொடக்க நிலை~$q_0$ இல் இருந்து, கட்டம்~$1$ ஐ வருடுகிறது:
\[
\Obj Q_{q_0}(\num{1}, \num{0})
\]
\item முதல் $k+1$ கட்டங்கள் $\TMendtape$,
  $\sigma_{i_1}$, \dots, $\sigma_{i_k}$ என்ற குறிகளைக் கொண்டுள்ளன:
\[
\Obj S_\TMendtape(\num{0}, \num{0}) \land
\Obj S_{\sigma_{i_1}}(\num{1}, \num{0}) \land
\dots \land
\Obj S_{\sigma_{i_k}}(\num{k}, \num{0})
\]
\item மற்றபடி, நாடா வெற்றாக உள்ளது:
\[
\lforall[x][(\num{k} < x \lif \Obj S_\TMblank(x, \num{0}))]
\]
\end{enumerate}
\item ஓர் அமைவுநிலையிலிருந்து அடுத்த அமைவுநிலைக்குச் செல்லும்
  நிலைமாற்றத்தை விவரிக்கும் அடிகோள்கள்:

பின்வருவனவற்றுக்காக, கீழுள்ள வடிவில் அமையும் எல்லா !!{sentence}s
இன் இணையலாக $!A(x, y)$ இருக்கட்டும்:
\[
\lforall[z][
  (((z < x \lor x < z) \land \Obj S_\sigma(z, y))
  \lif \Obj S_\sigma(z, y'))]
\]
இங்கு $\sigma \in \Sigma$. ``கட்டம்~$m$ ஐத் தவிர, $n+1$
படிகளுக்குப் பிறகுள்ள நாடா, $n$ படிகளுக்குப் பிறகுள்ள நாடாவைப்
போலவே உள்ளது'' என்பதை வெளிப்படுத்த $!A(\num{m},\num{n})$ ஐப்
பயன்படுத்துகிறோம்.
\begin{enumerate}
\item \ollabel{rep-right} ஒவ்வொரு கட்டளை $\delta(q_i, \sigma) =
  \tuple{q_j, \sigma', \TMright}$ க்கும், பின்வரும் !!{sentence}:
\begin{align*}
& \lforall[x][\lforall[y][(
   (\Obj Q_{q_i}(x, y) \land \Obj S_{\sigma}(x, y)) \lif {}]] \\
&\qquad   (\Obj Q_{q_j}(x', y') \land \Obj S_{\sigma'}(x, y') \land
!A(x, y)))
\end{align*}
இது பின்வருவதைக் கூறுகிறது:~$y$ படிகளுக்குப் பிறகு பொறி நிலை~$q_i$
இல் இருந்து, குறி~$\sigma$ ஐக் கொண்ட கட்டம்~$x$ ஐ வருடுகிறது எனில்,
$y+1$ படிகளுக்குப் பிறகு அது கட்டம்~$x+1$ ஐ வருடுகிறது, நிலை~$q_j$
இல் உள்ளது, கட்டம்~$x$ இப்போது~$\sigma'$ ஐக் கொண்டுள்ளது, மேலும்
~$x$ அல்லாத ஒவ்வொரு கட்டமும்~$y$ படிகளுக்குப் பிறகு கொண்டிருந்த அதே
குறியைக் கொண்டுள்ளது.

\item \ollabel{rep-left} ஒவ்வொரு கட்டளை $\delta(q_i, \sigma) =
  \tuple{q_j, \sigma', \TMleft}$ க்கும், பின்வரும் !!{sentence}:
\begin{align*}
& \lforall[x][\lforall[y][
    ((\Obj Q_{q_i}(x', y) \land \Obj S_{\sigma}(x', y)) \lif {}]]\\
& \qquad   (\Obj Q_{q_j}(x, y') \land \Obj S_{\sigma'}(x', y') \land
!A(x, y))) \land {}\\
& \lforall[y][((\Obj Q_{q_i}(\num{0}, y) \land \Obj S_{\sigma}(\num{0},
    y)) \lif {}]\\
& \qquad (\Obj Q_{q_j}(\num{0}, y') \land \Obj S_{\sigma'}(\num{0},
  y') \land !A(\num{0}, y)))
\end{align*}
இது எவ்வாறு செயல்படுகிறது என்று சற்று சிந்திக்க: இப்போது
``கட்டம்~$x$ ஐ வருடினால் \dots'' என்பதிலிருந்து தொடங்காமல்,
``கட்டம் $x+1$ ஐ வருடினால் \dots'' என்பதிலிருந்து தொடங்குகிறோம்.
இடப்புறம் நகர்வது என்பது அடுத்த படியில் பொறி கட்டம்~$x$ ஐ வருடுகிறது
என்று பொருள். ஆனால் எழுதப்படும் கட்டம்~$x+1$. நம்மிடம் கழித்தலோ
முன்னெண் சார்போ இல்லாததால் இதை இவ்வாறு செய்கிறோம்.

$x+1$ வடிவிலுள்ள எண்கள் $1$, $2$, \dots என்பதைக் கவனிக்க; அதாவது,
பொறி கட்டம்~$0$ ஐ வருடிக்கொண்டு இடப்புறம் நகர வேண்டிய நேர்வை இது
உள்ளடக்கவில்லை (இயல்பாகவே அதனால் அவ்வாறு நகர முடியாது---இருக்கும்
இடத்திலேயே அது தங்குகிறது). அச்சிறப்பு நேர்வை இரண்டாவது இணையல்
உள்ளடக்குகிறது: $y$ படிகளுக்குப் பிறகு பொறி நிலை~$q_i$ இல் இருந்து
கட்டம்~$0$ ஐ வருடுகிறது, மேலும் கட்டம்~$0$ குறி~$\sigma$ ஐக்
கொண்டுள்ளது எனில், $y+1$ படிகளுக்குப் பிறகும் அது கட்டம்~$0$ ஐ
வருடுகிறது, இப்போது நிலை~$q_j$ இல் உள்ளது, கட்டம்~$0$ இலுள்ள குறி
$\sigma'$, மேலும் கட்டம்~$0$ அல்லாத கட்டங்கள் $y$ படிகளுக்குப்
பிறகு கொண்டிருந்த அதே குறிகளைக் கொண்டுள்ளன என்று அது கூறுகிறது.
\item \ollabel{rep-stay} ஒவ்வொரு கட்டளை $\delta(q_i, \sigma) =
  \tuple{q_j, \sigma', \TMstay}$ க்கும், பின்வரும் !!{sentence}:
\begin{align*}
& \lforall[x][\lforall[y][(
   (\Obj Q_{q_i}(x, y) \land \Obj S_{\sigma}(x, y)) \lif {}]] \\
&\qquad   (\Obj Q_{q_j}(x, y') \land \Obj S_{\sigma'}(x, y') \land
!A(x, y)))
\end{align*}
\end{enumerate}
\end{enumerate}
டியூரிங் பொறி~$M$ மற்றும் உள்ளீடு~$w$ க்கான மேலுள்ள
!!{sentence}s அனைத்தின் இணையலாக $!T(M, w)$ இருக்கட்டும்.

~$M$ இறுதியில் நிற்கிறது என்பதை வெளிப்படுத்த, ``ஏதோ ஓர்
எண்ணிக்கையிலான படிகளுக்குப் பிறகு நிலைமாற்றுச் சார்பு
வரையறுக்கப்படாமல் இருக்கும்'' என்று கூறும் ஓர் !!{sentence} ஐக்
கண்டறிய வேண்டும்.~$\delta(q, \sigma)$ வரையறுக்கப்படாமல் உள்ள எல்லா
ஜோடிகள் $\tuple{q, \sigma}$ இன் கணமாக $X$ இருக்கட்டும். அப்போது
$!E(M, w)$ என்பது பின்வரும் !!{sentence} ஆக இருக்கட்டும்:
\[
\lexists[x][\lexists[y][(\bigvee_{\tuple{q, \sigma} \in
      X}(\Obj Q_q(x, y) \land \Obj S_\sigma(x, y)))]]
\]

ஒதுக்கப்பட்ட நிற்றல் நிலை~$h$ கொண்ட ஒரு டியூரிங் பொறியைப்
பயன்படுத்தினால் இது இன்னும் எளிது: அப்போது பின்வரும்
!!{sentence}~$!E(M, w)$
\[
\lexists[x][\lexists[y][\Obj Q_h(x, y)]]
\]
பொறி இறுதியில் நிற்கிறது என்பதை வெளிப்படுத்துகிறது.

\begin{prop}
\ollabel{prop:mlessk}
$m < k$ எனில், $!T(M, w) \Entails \num{m} < \num{k}$.
\end{prop}

\begin{proof}
பயிற்சி.
\end{proof}

\begin{prob}
\olref[tur][und][rep]{prop:mlessk} ஐ நிறுவுக.
(குறிப்பு: $k-m$ மீதான தொகுத்தறிதலைப் பயன்படுத்துக).
\end{prob}

\end{document}
